\chapter{用活神经元搭成的机器}
\chaptermark{用活神经元搭成的机器}
\label{sec:livingmachine}
\begin{chaptergoals}
\item 为什么没有广播通道的培养物会产生群放电，是什么决定其间隔；
\item 活的网络如何在反馈下改变响应，哪些仍未证实；
\item 类器官如何充当储备池，其线性读出在外部的硅片上训练；
\item 迄今培养物是怎样被教的，为什么老师仍站在培养皿之外。
\end{chaptergoals}

\section{电路图公开的试验台}

在第~\ref{sec:debugging}~章里，我们调试的是一台没有电路图的机器：探针在八百亿个神经元中只能听到一千个，其余全靠猜。处在这种境地的工程师会做一件简单的事：在面包板上搭一小块电路，每个元件都看得见。对神经元也可以这样做。

把皮层神经元分散开，培养在带电极阵列的芯片上，电极从几十个到几万个不等。几周后，细胞长出突起，彼此连接成由几千到几十万个神经元组成的网络，其中大部分是\nt{GLUT}，较小的一部分是\nt{GABA}，这一比例取决于标本：在海马培养物中约占神经元的 $6\%$~\cite{Benson1994}。每个电极既能像第~\ref{sec:debugging}~章的探针那样监听，也能像测试信号那样注入电流。试验台的第二种形式是\emph{类器官}：由人类干细胞培养出的三维神经组织小团，大小约一毫米~\cite{Lancaster2013}。在这样的试验台上，活的网络可以工作好几个月；有的平台允许通过网络远程在类器官上做实验，连续一百多天~\cite{Jordan2024}。这一领域被称为\emph{生物计算}或\emph{类器官智能}~\cite{Smirnova2023}。

试验台与大脑有两个重要区别。它极小：神经元数量少十万倍。而且它没有广播通道：皮层培养物中没有\nt{DOPA}、\nt{SERO}、\nt{NORA}和\nt{ACET}神经元。这是一台有数据总线和流量控制、却没有控制寄存器的机器。我们来看看它能做什么。

\section{网络自己在做什么}

放着不管的培养物既不沉默，也不均匀地产生噪声。每隔一段时间，几乎整个网络会同时爆发一次\emph{群放电}，然后又安静下来；群放电之间的间隔从一秒到几分钟不等，因培养物而异~\cite{Wagenaar2006}。对工程师来说这是熟悉的画面：带强正反馈、又没有负载的放大器会变成振荡器。

机制我们已经知道。\nt{GLUT}神经元之间强烈的相互连接会引发雪崩，就像第~\ref{sec:gaba}~章中命题~\ref{prop:ei} 的条件被破坏时那样。雪崩迅速耗尽突触中的神经递质储备~\eqref{eq:depression}，网络随之沉寂。储备以时间常数 $\tau_{\text{rec}}$ 恢复，当恢复到足以引发新雪崩的比例 $x^*$ 时，下一次群放电到来。群放电之间的间隔为
\begin{empheq}[box=\keybox]{equation}
T_{\text{burst}} \;\approx\; \tau_{\text{rec}}\,\ln\frac{1}{1-x^*} .
\label{eq:burstinterval}
\end{empheq}
取第~\ref{sec:code}~章的 $\tau_{\text{rec}}=0.8\,\text{s}$ 和 $x^*=0.9$，约为两秒；$x^*=0.99$ 时约为四秒。这是用本书 $\tau_{\text{rec}}$ 得到的模型估计，落在实测间隔的短端。在微培养物上的直接测量表明，群放电后递质储备要几十秒才能恢复~\cite{CohenSegal2011}，也就是说那里的 $\tau_{\text{rec}}$ 更大；用这样的 $\tau_{\text{rec}}$，公式~\eqref{eq:burstinterval} 给出几十秒乃至几分钟，与慢速培养物一致。这是一个弛张振荡器，与第~\ref{sec:muscles}~章的节律发生器同属一类：那里给它“充电”的是肌群的疲劳，这里是递质储备。

群放电是网络无事可做时所做的事。在活的大脑中，类似的图景出现在深度睡眠中（第~\ref{sec:sleep}~章），也出现在真正的输入尚未接入的发育期大脑中。这种相似性发人深思，但两者机制相同只是一个假说。

\section{没有多巴胺老师，怎样教网络}

培养物中没有\nt{DOPA}，“更好还是更差”的信号只能由我们自己通过电极给出。实验表明，试验台上的网络确实会改变响应，而其中最有意思的是用\emph{什么}来教它。

\emph{会闭嘴的老师。}通过一个电极每一到三秒刺激一次网络，直到另一个电极在刺激后 $50\pm10\ms$ 出现所需的响应。一旦出现，就关掉刺激五分钟。经过几个这样的循环后，所需响应在刺激后立刻出现~\cite{Shahaf2001}。这里的奖励是安静：刺激会扰动网络，而停止刺激则让网络停留在给出正确响应的那个状态。这就是第~\ref{sec:dofa}~章中通过随机偏离做梯度下降（命题~\ref{prop:gradient}），其中扮演误差角色的是“被扰动”这件事本身：一直扰动，直到对为止。

\emph{会打网球的网络。}在一项使用高密度电极阵列试验台的实验中，约一百万个神经元被接入一个单球拍的电脑网球游戏~\cite{Kagan2022}。球的位置以电流形式送到“感觉”电极上：球在哪里用位置编码，距离多远用频率编码。两个“运动”区的活动让球拍上下移动。学习规则很不寻常：如果球拍击中了球，网络得到一个短暂的\emph{可预测}刺激；如果没击中，则得到几秒钟\emph{不可预测}的随机电流。二十分钟的游戏中，连续击球的回合变长了。只有在完整反馈下，一次会话内才有显著提升；把刺激换成安静时，培养物到会话结束时也比没有反馈时打得好，但效果较小；而在几天里跨会话的稳定学习并未出现。对这项实验的详细分析见第~\ref{sec:dishbrain}~章。

作者用一个假说来解释这一结果：网络的行为是为了让自己的输入变得可预测~\cite{Friston2010}。这与第~\ref{sec:code}~章中节省脉冲、第~\ref{sec:recognition}~章中预测画面是同一个思想：输入在意料之中时，机器花费更少。但这项实验本身，尤其是作者关于培养物“有感知能力”的说法，招致了尖锐的批评：效应不大，而且可以用更简单的方式解释~\cite{Balci2023}。诚实的评价是：试验台上的培养物在反馈作用下改变行为，这是已测量的；至于它学的恰恰是预测自己的输入，则是假说。

\emph{会驾驶身体的网络。}以类似方式，人们把培养物接到一个简单的虚拟“身体”上，通过挑选训练刺激的时空模式，教它朝目标移动~\cite{Bakkum2008}。这些实验中都没有多巴胺：充当老师的是工程师选定的电流模式。当老师换成程序或多巴胺本身时会有什么变化，见第~\ref{sec:dishdopamine}--\ref{sec:cartpole}~节。

\begin{figure}[t]
\centering
\begin{tikzpicture}[>=Stealth,
  blk/.style={draw,very thick,rounded corners=2pt,align=center,font=\footnotesize,minimum height=0.8cm,fill=blue!5}]
% (a) closed loop
\node[font=\small] at (1.5,4.3) {(a) 闭环};
\node[blk,text width=2.6cm] (G) at (1.5,3.3) {游戏：球和球拍};
\draw[very thick,fill=orange!10,rounded corners=3pt] (0.5,0.2) rectangle (2.5,1.6);
\foreach \x in {0.7,1.0,...,2.35}{\foreach \y in {0.4,0.7,1.0,1.3}{\fill[gray!60] (\x,\y) circle (0.04);}}
\draw[->,thick,blue!60!black] (G.west) -| (-0.5,0.9) -- (0.5,0.9);
\node[font=\scriptsize,blue!60!black,anchor=east] at (-0.6,2.1) {球 $\to$ 电流};
\draw[->,thick,red!70!black] (2.5,0.9) -- (3.5,0.9) |- (G.east);
\node[font=\scriptsize,red!70!black,anchor=west] at (3.6,2.1) {脉冲 $\to$ 球拍};
\node[font=\scriptsize] at (1.5,-0.2) {电极阵列上的神经元};
\node[font=\scriptsize,align=center] at (1.5,-0.85) {击中：可预测刺激\\未击中：随机电流};
% (b) reservoir
\begin{scope}[xshift=7.9cm]
\node[font=\small] at (1.6,4.3) {(b) 储备池};
\node[blk,text width=1.1cm] (X) at (-0.4,1.0) {输入\\$u(t)$};
\draw[very thick,fill=orange!10] (1.6,1.0) circle (0.8);
\foreach \a/\r in {20/0.35,80/0.5,150/0.3,210/0.55,280/0.4,330/0.2,110/0.15}{\fill[red!60!black] ({1.6+\r*cos(\a)},{1.0+\r*sin(\a)}) circle (0.05);}
\node[font=\scriptsize] at (1.6,-0.2) {类器官：无监督};
\node[blk,text width=1.8cm,fill=green!8] (W) at (4.0,1.0) {读出\\$W_{\text{out}}$};
\draw[->,thick] (X) -- (0.8,1.0);
\draw[->,thick] (2.4,1.0) -- node[above,font=\scriptsize] {$\mathbf r(t)$} (W);
\draw[->,thick] (W) -- (4.0,2.3) node[above,font=\scriptsize] {$y(t)$};
\node[font=\scriptsize,green!40!black] at (4.0,0.3) {有监督训练};
\end{scope}
\end{tikzpicture}
\caption{让活网络做计算的两种方式。(a) 闭环：球的位置以电流输入，运动区的脉冲移动球拍，击球后的可预测刺激或随机刺激充当老师。(b) 储备池~\eqref{eq:reservoir}：不给类器官任何教学信号，它把输入分解成许多带记忆的非线性响应；按误差学习的只有外部的一个线性读出。类器官自身的连接同时也在变化，但是无监督的。}
\label{fig:livingmachine}
\end{figure}

\section{活的储备池}

使用活网络还有另一种方式：根本不去教它。在工程上这称为\emph{储备池计算}~\cite{Maass2002,JaegerHaas2004}。取一个现成的、带记忆的非线性系统，什么系统都行，只要它对输入的响应足够丰富，并且依赖于最近的过去。输入 $u(t)$ 激励它，得到响应向量 $\mathbf r(t)$，只训练线性读出
\begin{empheq}[box=\keybox]{equation}
y(t) \;=\; W_{\text{out}}\,\mathbf r(t) ,
\label{eq:reservoir}
\end{empheq}
其权重用第~\ref{sec:motorlearning}~章的最小二乘规则~\eqref{eq:lms} 求得。储备池做的正是第~\ref{sec:motorlearning}~章里小\nt{GLUT}神经元做的事：把输入展开成一个长向量，使所需的类别在其中可以用平面分开（第~\ref{sec:dendrites}~章）。

电极阵列上的类器官正好可以充当这样的储备池。把语音以电流模式输入给它，只训练一个线性读出后，辨别说话人的准确率约为 $78\%$~\cite{Cai2023}。$78\%$ 这个数字是作者自己报告、媒体跟着转述的。这个结果有用，但要弄清其中的“智能”在哪里：类器官提供非线性和衰减的记忆，按误差学习是在外部、在硅片上完成的（图~\ref{fig:livingmachine}）。同时，类器官本身也并非一成不变：据作者报告，训练期间它自身的功能连接发生了变化，他们称之为类器官内部的无监督学习~\cite{Cai2023}。准确率中有多少来自这种变化、多少来自读出，并未区分。

\section{试验台告诉我们关于这台机器的什么}

\emph{能量。}按估计~\eqref{eq:energy}，一个脉冲约耗 $10^{-10}$~J。一百万个神经元的培养物，每个神经元每秒发放一个脉冲，用于信号传递的功耗为
\begin{empheq}[box=\keybox]{equation}
P \;\approx\; 10^{6}\times1\,\text{s}^{-1}\times10^{-10}\,\text{J} \;=\; 10^{-4}\,\text{W},
\label{eq:dishpower}
\end{empheq}
即十分之一毫瓦。刺激、记录和处理这些脉冲的电子设备功耗要高出几个数量级。与第~\ref{sec:debugging}~章一样，瓶颈不在计算，而在与计算的连接。

\emph{缺少的部件。}试验台展示了带\nt{GABA}流量控制、却没有控制寄存器的\nt{GLUT}总线能做多少事：自组织、产生群放电、在反馈作用下改变响应，以及充当好的储备池。没有控制寄存器，它做不到的是自己决定什么算成功。在试验台上这个信号由工程师给出，而在大脑中，正如我们在第~\ref{sec:dofa}--\ref{sec:acet}~章所见，由广播通道给出。

\emph{伦理。}类器官由人类细胞培养而成，它们越复杂，这种组织能否有所感受的问题就越严肃。工程视角给出一个部分答案：第~\ref{sec:pain}~章中的疼痛不是某一个传感器的信号，而是一整套带闸门、音量调节器和广播通道的警报系统，而试验台上没有这些。但这只是推理，不是证明；该领域本身也主张从一开始就让伦理评估与实验同步进行~\cite{Smirnova2023}。

\begin{keyblock}
\begin{proposition}[试验台上的活网络]
\label{prop:livingmachine}
电极阵列上由\nt{GLUT}和\nt{GABA}神经元组成的网络会自发产生群放电~\eqref{eq:burstinterval}，在反馈作用下改变响应，并可充当储备池~\eqref{eq:reservoir}，供外部训练的读出使用。这些都是已测量的。至于这样的网络按某种普遍规则学习，例如学习预测自己的输入，则是假说；而关于其“有感知能力”的高调说法，大多数专家并不接受。
\end{proposition}
\end{keyblock}

\section{闪光释放的多巴胺}
\label{sec:dishdopamine}

据我们所知，唯一一项直接让多巴胺给培养物当老师的实验，是在一个研究者可远程接入的类器官平台上完成的~\cite{Jordan2024}。在浸泡类器官的液体中加入装在光敏“笼”里的多巴胺：被笼锁住的分子不作用于受体。笼锁多巴胺的浓度为 $1$~mM。需要奖励网络时，用紫外光照射：笼分解，多巴胺释放到周围空间。

回路是这样的。同一个刺激反复施加；如果它引发的脉冲比以前多，就打开闪光、释放多巴胺。在 $240$ 次重复中，诱发脉冲数总体上在增加。作者自己写道，仅凭这一项实验不能断定增长是多巴胺引起的~\cite{Jordan2024}。

对工程师来说，这是在培养皿中搭建三因子规则~\eqref{eq:threefactor} 的第一次尝试：发送者和接收者的活动留下资格迹，由一个公共信号决定是否把变化固定下来。但这里的信号只能为正：有奖励或没有奖励，装置给不出负误差 $\delta<0$。而且多巴胺扩散和清除都很慢，就像~\eqref{eq:lowpass} 中的浓度那样，因此频繁的闪光可能只是抬高了它的总体水平。要把学习与这种效应区分开，需要两个对照：在与网络响应无关的随机时刻给同样多的闪光；以及不加笼锁多巴胺、只给同样的闪光。还需要休息后的检验：多巴胺消散后，变化是否依然存在。

\section{多巴胺教硅片}
\label{sec:biohybrid}

在反方向的实验中，活细胞教电子器件~\cite{Keene2020}。把分泌多巴胺的细胞培养在一个有机神经形态元件上方的微通道中，这个元件是一只晶体管，其沟道电导充当突触权重。细胞释放的多巴胺在元件电极处被氧化，从而改变电导。该器件还模拟了突触间隙中递质的清除机制，因此权重不仅能长期改变，还能复原。

从电路角度看，这是一个化学输入的泄漏积分器：权重累积多巴胺流量，而“清除”决定它遗忘的快慢，与~\eqref{eq:lowpass} 是同一个RC电路。这里的关键是连接方向。第~\ref{sec:debugging}~章的探针看不到广播寄存器，因为它们是化学的。这个元件读取的恰恰是化学寄存器，并把它转换成电学权重。对脑机接口来说，这是通向一种传感器的路径：它不仅能听到数据总线，还能听到控制寄存器。

\section{自带多巴胺神经元的类器官}
\label{sec:midbrainorg}

人们已经能用人类干细胞培养出类器官，它们类似于大脑中\nt{DOPA}神经元所在的那个区域。其中有能工作、能分泌多巴胺的多巴胺神经元~\cite{Jo2016}。这类类器官主要用于研究疾病。我们没有找到用它们的多巴胺给另一个网络当老师的实验。

对于用活神经元搭成的机器，这正是缺失的部件。第~\ref{sec:livingmachine}~章的试验台是没有控制寄存器的数据总线加流量控制；这里已经有了广播信号源。缺的是评论器：一个计算预测误差~\eqref{eq:ctd} 并控制这些神经元的网络。把皮层类器官和这样的类器官连起来，让多巴胺按误差释放，是一个开放问题，目前只是假说，还不是实验。

\section{不靠多巴胺平衡杆子的类器官}
\label{sec:cartpole}

近期最完整的一项实验完全不用多巴胺~\cite{Robbins2026}。把小鼠皮层类器官接入“小车倒立摆”任务的闭环：类器官的活动推动小车，杆子的位置通过刺激反馈回来。两次尝试之间，类器官接受短暂的高频教学刺激。具体给哪些刺激，由一个强化学习程序选择。

结果是测量出来的：
\begin{itemize}
\item 对大多数类器官，程序选择的教学信号比随机选择的信号或没有信号效果更好；
\item 休息 $45$ 分钟后，改进没有保留；
\item 阻断两类谷氨酸受体AMPA和NMDA，改进消失。
\end{itemize}

最后一条说明学到的东西存在哪里：在\nt{GLUT}突触中，并且经由第~\ref{sec:weightstore}~节中充当输入—输出符合检测器的那个受体。第二条说明缺了什么：有快速变化，没有巩固，就像第~\ref{sec:motorlearning}~章中只有快学习者、没有慢学习者的情形。老师的角色则变了：挑选电流模式的工程师被程序取代，但老师仍然站在培养皿之外。

在更早的电极阵列培养物学习实验中，我们也没有找到任何一项使用多巴胺的：教学信号一律是电刺激。

\section{谁在教培养物}

\begin{table}[t]
\centering
\footnotesize
\setlength{\tabcolsep}{4pt}
\renewcommand{\arraystretch}{1.15}
\begin{tabular}{>{\raggedright}p{2.7cm}>{\raggedright}p{2.5cm}>{\raggedright}p{3.2cm}>{\raggedright\arraybackslash}p{4.4cm}}
\toprule
实验 & 谁来教 & 教学信号 & 已知情况 \\
\midrule
出现所需响应时停止刺激 & 工程师 & 刺激结束 & 响应被固定——已测量 \\
DishBrain（第~\ref{sec:dishbrain}~章） & 工程师 & 可预测电流或随机电流 & 一次会话内回合数的显著增长只出现在完整反馈下，安静条件下效应较小——已测量；跨会话不稳定；原因仍有争论 \\
小车倒立摆 & 强化学习程序 & 程序选择的高频刺激 & 优于随机刺激，但休息后不保留——已测量 \\
闪光释放多巴胺 & 规则“脉冲更多即奖励” & 光释放的多巴胺 & 脉冲增加——观察结果；多巴胺的作用未证实 \\
生物混合突触 & 活细胞 & 细胞分泌的多巴胺 & 元件权重的长期改变与复原——已测量 \\
含\nt{DOPA}神经元的类器官 & --- & --- & 有多巴胺来源；未用作老师 \\
\bottomrule
\end{tabular}
\caption{谁在用什么教芯片上的活神经元。}
\label{tab:dishteachers}
\end{table}

在所有由培养物本身学习的实验中，老师目前都在外面（表~\ref{tab:dishteachers}）：工程师、程序，或工程师设定的规则。多巴胺只进过培养皿一次，它的作用还有待证明。在培养皿中搭建真正的广播通道，像第~\ref{sec:dofa}~章那样带有评论器、误差和资格迹，是下一步，而这一步还没有人迈出。


\section{小结}

\begin{table}[t]
\centering
\footnotesize
\setlength{\tabcolsep}{4pt}
\renewcommand{\arraystretch}{1.15}
\begin{tabular}{>{\raggedright}p{2.8cm}>{\raggedright}p{4.2cm}>{\raggedright\arraybackslash}p{5.8cm}}
\toprule
实验 & 活网络做什么 & 已知情况 \\
\midrule
无输入的网络 & 群放电，间隔从一秒到几分钟~\eqref{eq:burstinterval} & 已测量；通过突触耗竭的机制是公认模型 \\
会闭嘴的老师 & 关掉刺激时，所需响应被固定下来 & 已测量 \\
打网球 & 连续击球回合变长；一次会话内的显著提升只出现在完整反馈下 & 行为变化已测量；跨会话学习不稳定；效应大小和“预测输入”的解释有争议 \\
虚拟身体 & 在挑选好的刺激模式下朝目标移动 & 已测量 \\
储备池 & 为读出~\eqref{eq:reservoir} 提供非线性和记忆 & 已测量；按误差学习在外部、在硅片上完成；类器官的连接也在变化 \\
能量 & 每一百万个神经元约 $10^{-4}$~W~\eqref{eq:dishpower} & 按~\eqref{eq:energy} 估计 \\
\bottomrule
\end{tabular}
\caption{用活神经元搭成的机器展示了什么。}
\label{tab:livingmachine}
\end{table}

用活神经元搭成的机器（表~\ref{tab:livingmachine}）是一个试验台，从中可以看出：没有控制寄存器时，数据总线和流量控制能做些什么。它们能做很多事，却不能自己决定什么是好。在试验台上，这个角色由工程师和他的电流模式扮演；在大脑中，由本书中间部分所讲的那几条广播导线扮演。

\section{习题}

\begin{exercise}[\lvlA]
\label{ex:21:1}
\emph{群放电间隔。}一次群放电把储备耗尽到零，随后 $\tau_{\text{rec}}\,\dd x/\dd t=1-x$，当 $x=x^*$ 时开始新的群放电；$\tau_{\text{rec}}=0.8$~s。(a)~求 $x^*=0.9$ 和 $0.99$ 时的间隔。(b)~阈值 $x^*$ 在 $0.99$ 附近随机变化 $\pm0.005$，间隔落在什么范围内？
\end{exercise}

\begin{exercise}[\lvlA]
\label{ex:21:2}
\emph{试验台的能量。}(a)~对整个大脑（$8.6\cdot10^{10}$ 个神经元），估计~\eqref{eq:dishpower} 给出多大功率？(b)~试验台以 $20$~kHz的采样率记录 $1024$ 个电极，每个采样 $10$~bit。按每bit $1$~nJ计，传输这一数据流的功耗是培养物本身功耗的多少倍？
\end{exercise}

\begin{exercise}[\lvlB]
\label{ex:21:3}
\emph{设计：抑制群放电。}通过许多电极的稀疏刺激使每个突触以频率 $r_b$ 释放递质；储备满足 $\tau_{\text{rec}}\,\dd x/\dd t=1-x-Ur_b\tau_{\text{rec}}x$，$U=0.5$，$\tau_{\text{rec}}=0.8$~s。$r_b$ 至少多大时，储备长不到 $x^*=0.9$；$0.99$？此时突触减弱多少？
\end{exercise}

\begin{exercise}[\lvlB]
\label{ex:21:4}
\emph{储备池的记忆不超过神经元数。}有 $N$ 个输出 $\mathbf r(t)$ 的储备池接收零均值、单位方差的白噪声输入 $u(t)$；$m_k$ 是读出 $\mathbf w_k\cdot\mathbf r(t)$ 与 $u(t-k)$ 的相关系数平方的最大值。证明记忆容量 $\mathrm{MC}=\sum_{k\ge0}m_k\le N$。
\end{exercise}

\begin{exercise}[\lvlB]
\label{ex:21:5}
\emph{建模：硅片中的储备池。}储备池 $\mathbf r(t+1)=\tanh(W\mathbf r(t)+\mathbf v\,u(t)+\mathbf c)$，$N=30$，$W$ 为随机矩阵，谱半径 $0.9$，$v_i\sim U(-0.5,0.5)$，$c_i\sim U(-0.2,0.2)$，$u\sim U(-1,1)$，读出用岭回归。求带 $\tanh$ 和不带它时习题~\ref{ex:21:4} 的记忆容量。哪个储备池记得更多？
\end{exercise}

\begin{exercise}[\lvlB]
\label{ex:21:6}
\emph{为活储备池设计读出。}类器官提供 $1024$ 个通道和两类共 $P=240$ 条训练记录。(a)~按习题~\ref{ex:19:3} 的公式，最多保留多少个特征 $n$，才能使随机标注可分的概率不超过 $1\%$？(b)~在 $100$ 条测试记录上 $78\%$ 的准确率比随机猜测高出几个标准差？
\end{exercise}

\begin{exercise}[\lvlC]
\label{ex:21:7}
\emph{研究：没有广播通道的老师。}提出一个刺激方案，通过 $8$--$1024$ 个电极在试验台上实现第~\ref{sec:dofa}~章的三因子规则；再设计一个冻结读出的对照，用来区分权重学习与网络状态漂移。什么结果会否定“学习”这一假说？
\end{exercise}
